\section{Introduction}
\label{sec:intr}

\subsection{$q$-Web categories}

Two monoidal categories, the affine web category and the affine Schur category, were introduced by the authors \cite{SW25web, SW24Schur}, and their cyclotomic quotients were also constructed; also see \cite{DKM25} for a generalization of the affine and cyclotomic web categories. Two strict $\kk$-linear monoidal categories, the $q$-deformed affine web and Schur categories, $\qAW$ and $\qASch$, were subsequently formulated in \cite{SSW25}, and we will assume $\kk=\C[q,q^{-1}]$ in this paper. The thin strand full subcategory of $\qAW$ is identified with the affine Hecke algebra category $\qAH$.

These three monoidal categories and their corresponding cyclotomic quotient categories can be assembled into the following commutative diagram:
\begin{equation} \label{CD:categories}
\begin{tikzcd}
\qAH \ar[r,hook]\ar[d,two heads] & \qAW\ar[r,hook]\ar[d,two heads] & \qASch\ar[d,two heads]\\
\qH_\bfu  \ar[r,hook]&\qW_\bfu \ar[r,hook]&\qSch_\bfu
\end{tikzcd} 
\end{equation}
where $\bfu=(u_1,\ldots, u_\ell) \in (\kk^*)^\ell$, for $\ell \in\N$. These categories have provided (often new) diagrammatic presentations of old and new quantum algebras. In this paper, we are mostly focused on the algebras arising from the middle column and their relations to those from the left column.

The path algebra (i.e., the direct sum of all hom-spaces) of $\qAW$ is 
\[
\hwS:=\bigoplus_{n\ge 0} \hwS(n), 
\]
where it follows by \cite{SSW25} that $\hwS(n)$ is Morita equivalent to the affine Schur algebras $\ASch(n,n)$ from \cite{Gre99}.
%, while the path algebra of $\qSch$ is Morita equivalent to the higher level affine Schur algebras from \cite{MS21}. 
%(The Morita equivalences can be upgraded to isomorphisms if we adopt variants of these categories allowing compositions instead of strict compositions as objects.)
The path algebra of $\qW_\bfu$ is 
\[
\wS_\bfu:=\bigoplus_{n\ge 0} \wS_\bfu(n); 
\]
we propose to call these new algebras $\wS_\bfu(n)$ {\em cyclotomic $q$-W-Schur algebras} (of level $\ell$ if $\bfu$ is an $\ell$-tuple); recall that the degenerate variants were identified in \cite{SW25web} with Schur algebras for finite W-algebras. 
%We have inclusions of algebras $\bfH_\bfu(n) \subset \wS_\bfu(n) \subset \bfS_\bfu(n)$, for general $\bfu$. 
The level one cyclotomic $q$-web category appeared in \cite{Bru25} (called $q$-$\mathbf{Schur}$ category) and these 2 level one cyclotomic $q$-(W-)Schur algebras are identical and Morita equivalent to Dipper-James $q$-Schur algebras $\bfS(n,n)$. 

We also recall the well-known fact that the path algebra of $\qAH$ is identified with the affine Hecke algebras: $\AHe :=\oplus_{n\ge 0} \AHe(n)$, while the path algebra of the cyclotomic category $\qH_\bfu$ is identified with the (Ariki-Koike) cyclotomic Hecke algebras: $\bfH_\bfu =\oplus_{n\ge 0} \bfH_\bfu(n)$.




\subsection{Works of \cite{LLT96, Ar96, LT96, VV99} }

Let $\bfs=(s_1,\ldots, s_\ell)\in \Z^\ell$. Let $\eps \in \C$ be generic or a root of $1$, and let $p$ be the order of $\eps^2$. Denote by $\bfH_{\eps^{2\bfs}}(n)$ the cyclotomic Hecke algebras by specializing $q=\eps$ and $\bfu=\eps^{2\bfs}$. We denote by $[\mathcal C]$ the Grothendieck group of an abelian category $\mathcal C$. Denote by $V(\omega_\bfs)$ the integrable highest weight module $V(\omega_\bfs)$ over the affine quantum group $\U_v(\widehat\sll_p)$, and denote by $\varpi_\bfs: \U^-_v(\widehat\sll_p) \rightarrow V(\omega_\bfs)$ the $\U^-_v(\widehat\sll_p)$-module homomorphism sending $1$ to the highest weight vector $\eta_{\omega_\bfs}$. The surjection $\qH \rightarrow \qH_{\eps^{2\bfs}}$ induces a natural functor $\pi_\bfs: \qH_{\eps^{2\bfs}}\modf \rightarrow \qH\modf$ which preserves the simples, which in turn induces a surjective map on the duals, $\pi_\bfs^*: [\qH\modf]^* \rightarrow [\qH_{\eps^{2\bfs}}\modf]^*$. Proving and extending Lascoux-Leclerc-Thibon conjecture \cite{LLT96}, Ariki \cite{Ar96} established the following commutative diagram of $\U^-(\widehat\sll_p)$-module homomorphisms: 
\begin{equation}  
\label{eq:Ariki}
\begin{tikzcd}
\U^-_\Z(\widehat\sll_p) \ar[r,"\cong"]\ar[d,two heads,"\varpi_\bfs"]&{[}\qH\modf_\eps{]}^* \ar[d,two heads,"\pi_\bfs^*"]\\
V(\omega_\bfs)_\Z  \ar[r,"\cong"]&{[} \qH_{\eps^{2\bfs}}\modf{]}^* 
\end{tikzcd} 
\end{equation}
where the horizontal isomorphisms send the canonical basis \cite{Lus90, Lus91} (and \cite{Kas91}) specialized at $v=1$ to the dual basis of simple modules, and the Chevalley generators $F_i, E_i$ of $\U_v(\widehat\sll_p)$ are realized as $i$-restriction and $i$-induction functors on $\qH_{\eps^{2\bfs}}\modf$. Here and below we use the subscript $\Z$ to denote the specialization at $v=1$. We refer to the survey of Kleshchev \cite{Kle10} for numerous subsequent exciting developments on representations of symmetric groups, Hecke algebras, and categorification.

Varagnolo-Vasserot \cite{VV99} showed that the decomposition matrices of the $q$-Schur algebras for $q=\eps$ being a root of 1 are given by the transition matrices between the canonical basis and Schur basis on the level one $\U_v(\widehat\gl_p)$-Fock space $\cF$, establishing the Leclerc-Thibon conjecture \cite{LT96}; this generalizes the level one result in \cite{LLT96, Ar96}. Their results can be paraphrased as a $\Z$-linear isomorphism:
\begin{align}  \label{eq:VV}
     \cF_\Z \stackrel{\cong}{\longrightarrow} \bigoplus_{n\ge 0} [\bfS(n,n)\modf]^*,
\end{align}
which maps the Schur basis and the canonical basis to the dual bases to Weyl modules and simple modules, respectively. 

\subsection{The main results}

This paper is a first step on categorification based on our diagrammatic web and Schur categories as outlined in \cite{SW25web, SW24Schur}. In the setting of the middle column of the diagram \eqref{CD:categories}, we show that the cyclotomic $\eps$-webs with $\eps^2$ being a primitive $p$th root of $1$ categorify the integrable highest weight modules over quantum affine $\gl_p$, which admit canonical bases exhibiting positivity (see \cite{Sch00}); this generalizes the main results of Ariki and Varagnolo-Vasserot. Similar results hold in the degenerate web setting. To that end, it is crucial and fruitful for us to develop new connections of affine $q$-webs to geometric convolution algebras and to Hall algebra of the cyclic quiver. 

The results of this paper will be a building block of the categorification of additional diagrammatic web/Schur categories of other types. 

\subsubsection{}

One advantage of diagrammatic categories is to allow us to construct restriction and induction functors on $\qW_{\bfu} \modf$, which naturally interact with combinatorics of semi-standard tableaux. It is shown in \cite{SSW25} (see \cite{SW24Schur}) that the path algebra of $\qSch_\bfu$ is isomorphic to (Dipper-James-Mathas) cyclotomic $q$-Schur algebras \cite{DJM98}: $\bfS_\bfu=\oplus_{n\ge 0} \bfS_\bfu(n)$. The three families of algebras arising from the diagram \eqref{CD:categories} are related: $\bfH_\bfu(n)\subset \wS_\bfu(n) \subset \bfS_\bfu(n)$. We obtain the cellular structure on $\wS_\bfu(n)$ by applying a Schur functor to $\bfS_\bfu(n)\modf$; note that $\bfS_\bfu(n)$ is quasi-hereditary \cite{SSW25} and provides a diagrammatic realization of the cyclotomic $q$-Schur algebras \cite{DJM98}. Denote by $\Par^\ell(n)$ the set of $\ell$-multipartitions of $n$.


\begin{alphatheorem} [Theorem \ref{thm:cellularforwebu}, Proposition \ref{prop:IndRes}] 
\label{thm:A}
\begin{enumerate}
    \item The algebra $\wS_\bfu(n)$ admits a cellular basis and cell modules $\Delta^\la$, for $\la \in \Par^\ell(n)$.
    \item 
    There exist cell module filtrations on $\Res^n_{n-a}(\Delta^\la)$ and $\Ind_n^{n+a} (\Delta^\la)$, for $\la \in \Par^\ell(n)$. 
\end{enumerate}
\end{alphatheorem}

\subsubsection{}
Fix $\I =\Z$ or $\Z/p\Z$ when $q=\eps$ is generic or a root of 1 with order of $\eps^2$ equal to $p$, and choose the parameter $\bfu =\eps^{2\bfs}$. The quotient functor $\qAW \to\qW_{\eps^{2\bfs}}$ induces a functor 
\[
\pi_\bfs: \qW_{\eps^{2\bfs}}\modf \longrightarrow \qAW\modf_\eps
\]
and then a surjection $\pi_\bfs^*: [\qAW\modf_\eps]^* \rightarrow [\qW_{\eps^{2\bfs}}\modf]^*$. Denote by $\bfm\in\cM\cS$ the set of all multi-segments \eqref{cMcS} over $\I$. By definition \eqref{eq:std}, the standard modules $\Delta_\bfm$, for $\bfm\in\cM\cS$, are $\hwS(n)$-modules of the form $\pi_\bfs(\Delta^\la)$, for some distinguished $\la \in \Par^\ell(n)$ associated to $\bfm$. 

The geometric affine $q$-Schur algebra, $\KS(N,n) =K^G(T^*\cF \times_{\cN} T^*\cF)$, is the convolution algebra on the Steinberg variety of $N$-step flags in $\C^n$; cf. \eqref{KSchur}. Denote by 
\[
{\K} :=\bigoplus_{n=0}^\infty [\KS(n,n)\modf]
\]
and by $\K^*$ its restricted dual; cf. \cite{GV93, V98} or \eqref{eq:KK}. Denote by $M_\bfm$, for $\bfm\in\cM\cS$,  the standard modules in $\K$ defined via Borel-Moore homology; cf. \cite{CG97, VV99}. The theorem below can be viewed as a generalization of the induction theorem of Kazhdan-Lusztig \cite{KL87} for affine Hecke algebras as adapted by Ariki \cite{Ar96}.

\begin{alphatheorem} [Theorem \ref{thm:induced}, Proposition \ref{prop:Morita} and Theorem \ref{thm:ind=std}]
\label{thm:B}
\qquad
    \begin{enumerate}
        \item 
        For $\la \in \Par^\ell(n)$, the $\hwS(n)$-module $\pi_\bfs (\Delta^\la)$ admits an induced module structure. 
        \item The algebras $\hwS(n)$ and $\KS(n,n)$ are Morita equivalent, inducing an isomorphism $\natural: [\qAW\modf_\eps] \stackrel{\cong}{\longrightarrow}\K$.
        \item $\natural$ matches the standard modules:  $\natural([\Delta_\bfm] )= [M_\bfm]$, for all $\bfm\in\cM\cS$. In particular, $\{[{\Delta}_\bfm]\mid \bfm\in\cM\cS\}$ forms a $\Z$-linear basis for $[\qAW\modf_\eps].$
    \end{enumerate}
\end{alphatheorem}

%Part (1) of Theorem \ref{thm:B} is proved as Theorem \ref{thm:induced}. 
There is another version of affine $q$-Schur algebra $\ASch(N,n)$ defined in \eqref{def:affineSchur} as an endomorphism algebra \cite{Gre99}. It is known that (see \cite{LXY26}) 
\[
\KS(N,n) \cong \ASch(N,n)
\]
for $N\ge n$, and all $\ASch(N,n)$ for $N\ge n$ are Morita equivalent to $\ASch(n,n)$. On the other hand, it was shown in \cite{SSW25} that $\ASch(n,n)$ is Morita equivalent to $\hwS(n)$, whence Part (2). Part (3) follows from the induction theorem for the affine Hecke algebra \cite{KL87, Ar96} by applying the Schur functor and Part (1) (=Theorem \ref{thm:induced}). 

\subsubsection{}
Lusztig \cite{Lus91, Lus98} used the composition subalgebra of the generic Hall algebra $\Hall$ of a cyclic quiver $\pen_p$ and its IC basis to realize $\U^-_v(\widehat\sll_p)$ and its canonical basis. The generic Hall algebra $\Hall$ is isomorphic to half the affine quantum group $\U^-_\cA(\widehat\gl_p)$ with $\cA=\Z[v,v^{-1}]$, and its IC basis is regarded as the canonical basis of the latter; cf. \cite{LT96, VV99}. The Drinfeld double of $\Hall$ (i.e., double Hall algebra) over $\C(v)$ can be naturally identified with $\U_v(\widehat\gl_p)$; cf. \cite{X97, DF15, DX17}. 

In \cite[\S12.3]{VV99} (building on \cite{CG97, GV93, GRV94}), it was pointed out that there exists a $\Z$-linear isomorphism (see Theorem \ref{thm:Hall=Schursimple}):
\begin{align}  \label{Hall=K}
    \imath_p: \HallZ \stackrel{\cong}{\longrightarrow} \K^*,
\end{align}
which sends a Hall basis $f_\bfm$ (corresponding to characteristic functions on orbits in the Hall algebra over finite fields) to the basis of standard modules $[M_\bfm]$, for $\bfm\in\cM\cS$. Moreover, $\imath_p$ sends the canonical basis to the basis dual to the simple modules. 

The Hall algebra $\Hall$ is generated by the semisimple generators (corresponding to the semisimple quiver representations), and a closed multiplication formula for $f_\bfi f_\bfm$ between a semisimple generator $f_\bfi$ (for $\bfi\in\N\I$) and any Hall basis element is obtained by Du-Fu \cite{DF15}; also see \cite{FL19} for a geometric proof. We view $\Hall$ as a left regular $\U_\cA^-(\widehat\gl_p)$-module via left multiplication and the identification $\U^-_\cA(\widehat\gl_p)\cong\Hall$. 

On the other hand, we define the $\bfi$-restriction and $\bfi$-induction functors diagrammatically which give rise to the linear operators $f_\bfi$ on $[\qAW \modf_\eps]^*$ and $f_\bfi$ and $e_\bfi$ on $[\qW_{\eps^{2\bfs}} \modf]^*$. 

\begin{alphatheorem}
     [Theorem \ref{thm:Cat_affine}]
     \label{thm:C}
 Let $p\le \infty$.
 \begin{enumerate}
 \item The $f_\bfi$ via $\bfi$-restriction functors provide an action of $\U^-_\Z(\widehat\gl_p)$ on $[\qAW\modf_\eps]^*.$
     \item 
     There exists a $\U^-_\Z(\widehat\gl_p)$-module isomorphism
\[
\jmath_p: \HallZ \longrightarrow [\qAW\modf_\eps]^*, 
\qquad \jmath_p(f_\bfm) =[{\Delta}_\bfm]^*.
\]
\item 
The map $\jmath_p$ sends the canonical basis to the basis dual to the simple modules.
 \end{enumerate} 
\end{alphatheorem}

Theorem \ref{thm:C} (1)-(2) strengthens the basis-preserving linear isomorphism $\imath_p$ of Ginzburg, Varagnolo and Vasserot in \eqref{Hall=K}. We compute $\iRes$ on ${\Delta}_\bfm$ via Theorem \ref{thm:A} and show that it matches exactly with the Du-Fu multiplication formula on $\Hall$, whence (1)--(2). Part (3) follows from the following commutative diagram of $\Z$-linear isomorphisms (as there is no known geometric $\U^-_\Z(\widehat\gl_p)$-action on $\K^*$): 
\begin{equation}   \label{eq:ijk}
\begin{tikzcd}
\HallZ\ar[r,"\imath_p"]\ar[dr,"\jmath_p"] &\K^*\ar[d,"\natural^*"]
\\
 & {[}\qAW\modf_\eps{]}^*
\end{tikzcd} 
\end{equation}


\subsubsection{}
Denote by $\Lambda(\omega_\bfs)$ the integrable  $\U_v(\widehat\gl_p)$-module of highest weight $\omega_\bfs$. Denote by $\varpi_\bfs: \Hall \rightarrow \Lambda(\omega_\bfs)\!_\cA$ the $\U^-_\cA(\widehat\gl_p)$-module homomorphism sending $1$ to the highest weight vector $v^+_{\omega_\bfs}$. It was proved by Schiffmann \cite{Sch00} that the canonical basis on $\Hall$ descends via 
$\varpi_\bfs$ to a canonical basis on $\Lambda (\omega_\bfs)$, the latter being part of Uglov's canonical basis on a higher level Fock space \cite{Ugl00}; moreover, the canonical basis on $\Lambda (\omega_\bfs)$ admits a positive integral expansion into the Schur basis of the Fock space. This generalizes the level one conjecture in \cite{VV99} building on \cite{LT96}. The following provides a generalization of the main results of Ariki and Varagnolo-Vasserot.

\begin{alphatheorem}  [Proposition \ref{prop:lieaction}, Theorem \ref{thm:cyc_cat}]
    \label{thm:D}  
We have the following commutative diagram of $\U^-_\Z(\widehat\gl_p)$-module homomorphisms: 
\begin{equation}  
\label{diag:comm}
\begin{tikzcd}
{\HallZ} \ar[r,"\jmath_p"]\ar[d,two heads,"\varpi_\bfs"]&{[}\qAW \modf_\eps{]}^* \ar[d,two heads,"\pi_\bfs^*"]\\
\Lambda (\omega_\bfs)_\Z  \ar[r,"\jmath_\bfs"]&{[}\qW_{\eps^{2\bfs}} \modf{]}^* 
\end{tikzcd} 
\end{equation}
Moreover, the map $\jmath_\bfs$ is a $\U_\Z(\widehat\gl_p)$-module isomorphism and sends the canonical basis to the basis dual to simple modules. 
\end{alphatheorem}
Theorem \ref{thm:D} is consistent with (and recovers via Schur functors) the classic results of Ariki \cite{Ar96}; see the diagram \eqref{CD:HeckeWeb}.

Denote by $\qW_{\eps^{2\bfs}} \prmod$ the category of finite-dimensional projective $\qW_{\eps^{2\bfs}}$-modules. Via the Cartan pairing, one can replace $[\qW_{\eps^{2\bfs}}\modf]^*$ and the duals of the simple modules in Theorem \ref{thm:D}  by $[\qW_{\eps^{2\bfs}}\prmod]$ and indecomposable projective modules $P^\la$. Then Theorem \ref{thm:D} can be rephrased as computing the cell filtration multiplicities of $P^\la$ in terms of $v$-decomposition matrix (specialized at $v=1$) from  canonical basis to Schur basis as well as providing a classification of simple modules over the cyclotomic  $q$-$W$-Schur algebras; see Corollaries \ref{cor:PIM}--\ref{cor:classifysimple}.


%It suffices to prove Theorem \ref{thm:D} by a base change to $\C$. Note that $[\qW_{\eps^{2\bfs}} \modf]_\C^*$ is a cyclic $\U(\widehat\gl_p)$-module; moreover it is integrable as one can embed it into a higher level Fock space and the latter can be identified with the Grothendieck group of $\qW_{t^\bullet\eps^{2\bfs}} \modf$ where $t^\bullet\eps^{2\bfs}$ is a semisimple parameter. Theorem \ref{thm:D} follows from these observations and applying Theorem \ref{thm:C}. Theorem \ref{thm:D} when specialized at level one strengthens  Varagnolo-Vasserot's basis-preserving isomorphism \eqref{eq:VV}. 

\subsubsection{}
The clockwise orientation of a cyclic quiver (i.e. $i\mapsto i+1$, for $i\in\I$) is adopted in \cite{VV99, Sch00, DF15}, while the opposite orientation is used in \cite{Ar96, LTV99}. In this paper, we follow the clockwise convention mainly because we need to match our computation in $\qAW$ exactly with (the $v=1$ specialization of) the multiplication formula in $\Hall$ with semisimple generators from \cite{DF15}; see Remarks \ref{rem:sign}--\ref{rem:opposite}. 

The constructions of $\bfi$-restriction and $\bfi$-induction functors for cyclotomic $q$-web categories formulated in this paper work equally well for the cyclotomic $q$-Schur categories $\qSch_{\eps^{2\bfs}}$, which categorify Uglov's Fock space \cite{Ugl00}. This follows from 
\cite{RSVV} thanks to the identification of the path algebra of cyclotomic $q$-Schur categories with cyclotomic $q$-Schur algebras. Our diagrammatic constructions of $\bfi$-restriction and $\bfi$-induction functors are more general and conceptual than Wada's \cite{Wad14}.


\subsubsection{}
The paper is organized as follows. In Section \ref{sec:web}, we introduce the restriction and induction functors for the $q$-web categories. In Section \ref{sec:cellular}, we construct and study the cellular modules $\Delta^\la$ of $\wS_\bfu(n)$. A semisimplicity criterion for $\wS_\bfu(n)$ is given. In Section~ \ref{sec:standard}, we define the standard modules of $\hwS(n)$ as the pullback of cell $\wS_\bfu(n)$-modules, and show that they are always induced modules. Theorem \ref{thm:C} is the key categorification result of this paper and will be established in Section \ref{sec:affineSchur_Hall}. Theorem \ref{thm:D} is to a large extent a consequence of Theorem \ref{thm:C} and will be proved in Section \ref{sec:Cat}. We end the paper with a discussion of some open problems in \S\ref{ssec:open}. 


\vspace{2mm}
%\noindent {\bf Funding and Competing Interests.}
%
{\bf Acknowledgement.} 
We thank Ming Lu and Changjian Su for helpful discussions on Hall algebra and geometric representation theory. We thank Department of Mathematics and IMS at University of Virginia and at National University of Singapore, Academia Sinica and the NSTC in Taipei, as well as University of Hong Kong for support and facilitating this collaboration. L.S. is partially supported by NSFC (Grant No. 12071346). W.W. is partially supported by DMS--2401351. 

%The authors have no competing interests to declare that are relevant to the content of this article. The authors declare that the data supporting the findings of this study are available within the paper.